\bookchapter{บทที่ 41: ภาพรวมเครื่องมือเขียนโค้ดด้วย AI ในปี 2026}

\chapterintro{ขอเตือนล่วงหน้าก่อน: บทนี้เป็น “ภาพหยุดนิ่ง” ของระบบนิเวศเครื่องมือเมื่อเดือนสิงหาคม 2026 อีกครึ่งปีก็จะเก่า — ซึ่งก็คือบทเรียนแรกที่บทนี้อยากสอนคุณพอดี เป้าหมายของบทนี้ไม่ใช่สอนว่า “ต้องใช้เครื่องมือไหน” แต่ช่วยให้คุณสร้างกรอบการตัดสินใจว่า “จะเลือกเครื่องมืออย่างไร” พอคุณเข้าใจปรัชญาการออกแบบของเครื่องมือแต่ละแบบ คุณก็จะเลือกสิ่งที่เหมาะกับตัวเองได้ไม่ว่าเมื่อไหร่}

\booksection{สามสำนักและกรอบการเลือก}

เครื่องมือเขียนโค้ดด้วย AI ในปี 2026 แบ่งคร่าว ๆ ได้เป็นสามสำนัก แต่ละสำนักเบื้องหลังคือปรัชญาการออกแบบที่ต่างกัน ลองดูภาพรวมนี้ก่อน:

\begin{center}
\begin{tabularx}{\linewidth}{|X|X|X|X|}
\hline
\textbf{สำนัก} & \textbf{เครื่องมือตัวแทน} & \textbf{ปรัชญาการออกแบบ} & \textbf{สถานการณ์ที่เหมาะ} \tabularnewline
\hline
AI ในตัว IDE & Cursor 3, Trae, Windsurf, CodeBuddy & รวมความสามารถ AI เข้ากับตัวแก้ไขโค้ดอย่างลึกซึ้ง & พัฒนาประจำวัน รีวิวโค้ด เสริมความสามารถแบบค่อยเป็นค่อยไป \tabularnewline
Agent บนเทอร์มินัล & Claude Code, Codex CLI, Aider & ใช้เทอร์มินัลเป็นหน้าจอหลัก AI จัดการไฟล์และคำสั่งโดยตรง & งานระดับโปรเจกต์ แก้ไขเป็นชุด การผสาน CI/CD \tabularnewline
Builder บนคลาวด์ & v0, Bolt, Replit Agent & สร้างโปรเจกต์เต็มรูปแบบจากศูนย์ในเบราว์เซอร์ & ทำต้นแบบเร็ว เริ่มใช้ได้ทันทีไม่ต้องตั้งค่า พรีวิวได้ทันที \tabularnewline
\hline
\end{tabularx}
\end{center}

\begin{center}\small\color{gray}สามสำนักของเครื่องมือเขียนโค้ดด้วย AI\end{center}

กุญแจของการเลือกเครื่องมือไม่ใช่ “ตัวไหนดีที่สุด” แต่คือ “ตัวไหนเหมาะกับสถานการณ์ปัจจุบันของคุณที่สุด” มุมมองตัดสินใจที่ใช้ได้จริงสามด้าน:

\begin{itemize}
  \item พื้นฐานการเขียนโค้ดของคุณเป็นอย่างไร? เริ่มจากศูนย์เลือก Builder บนคลาวด์ มีพื้นฐานเลือก AI ในตัว IDE ผู้ใช้ระดับสูงเลือก Agent บนเทอร์มินัล
  \item โปรเจกต์ของคุณอยู่ช่วงไหน? เริ่มจากศูนย์เลือก Builder บนคลาวด์ พัฒนาไล่เวอร์ชันเลือก AI ในตัว IDE ดูแลรักษากับ refactor เลือก Agent บนเทอร์มินัล
  \item ความต้องการการ deploy ของคุณคืออะไร? ตรวจแนวคิดเร็ว ๆ เลือก Builder บนคลาวด์ โปรเจกต์จริงจังเลือก AI ในตัว IDE ทำ workflow อัตโนมัติเลือก Agent บนเทอร์มินัล
\end{itemize}

\booksection{สำนัก AI ในตัว IDE}

AI ในตัว IDE เป็นสำนักที่ได้รับความนิยมหลักที่สุด: ฝังความสามารถ AI ลงในตัวแก้ไขโค้ดของคุณ คงประสบการณ์เขียนโค้ดแบบดั้งเดิมไว้ พร้อมเสริมพลัง AI ทับอีกชั้น ในปี 2026 ตัวแทนของสำนักนี้คือ Cursor 3, Trae, Windsurf และ CodeBuddy

\booksubsection{Cursor 3}

Cursor เป็น AI IDE ที่คนใช้แพร่หลายที่สุดในโลกตอนนี้ Cursor 3 ที่เปิดตัวเมื่อเมษายน 2026 ถือเป็นจุดเปลี่ยนสำคัญของเครื่องมือเขียนโค้ดด้วย AI: อัตราการใช้โหมด Agent แซงโหมดเติมโค้ดด้วย Tab เป็นครั้งแรก — นักพัฒนาเปลี่ยนจาก “AI ช่วยเติมโค้ดให้ผม” มาเป็น “AI ช่วยทำงานให้ผมเสร็จ” ฟีเจอร์หลักของ Cursor 3 ได้แก่: แก้ไขหลายไฟล์พร้อมกัน (AI แก้หลายไฟล์พร้อมกันเพื่อทำงานให้เสร็จ), การรู้บริบท (เข้าใจโครงสร้างโปรเจกต์และความสัมพันธ์ระหว่างไฟล์อัตโนมัติ) และการผสานกับเทอร์มินัลอย่างลึกซึ้ง (AI รันคำสั่งและดูผลลัพธ์ได้จากในตัวแก้ไขโค้ดเลย)

\booksubsection{Trae}

Trae เป็น AI IDE จาก ByteDance มีฐานผู้ใช้ขนาดใหญ่ในตลาดจีน โหมด SOLO (หรือเรียกว่าโหมด Builder) ของมันคือความสามารถ “พูดหนึ่งประโยคสร้างโปรเจกต์” อย่างแท้จริง — คุณบอกความต้องการ AI จัดการครบทั้งกระบวนการตั้งแต่เริ่มโปรเจกต์จนถึงสร้างโค้ดด้วยตัวเอง จุดแข็งของ Trae คือรองรับภาษาจีนและปรับให้เข้ากับสภาพแวดล้อมเครือข่ายในประเทศ เหมาะกับผู้ใช้ในจีนที่อยากเริ่มใช้งานได้เร็ว

\booksubsection{เครื่องมืออื่น ๆ}

Windsurf (เดิมคือ Codeium) ขึ้นชื่อเรื่องประสบการณ์เติมโค้ดที่ลื่นไหล เหมาะกับนักพัฒนาที่อยากได้ความรู้สึก “พิมพ์เร็วปานสายฟ้า” CodeBuddy ของ Tencent Cloud เป็นปลั๊กอินของ VS Code เก่งเรื่องผสานกับระบบนิเวศคลาวด์ เหมาะกับโปรเจกต์ที่วางแผน deploy บน Tencent Cloud หรือพัฒนา WeChat Mini Program ส่วน Tongyi Lingma ของ Alibaba และ Comate ของ Baidu ก็มอบประสบการณ์ผสานกับ IDE แบบคล้ายกัน โดยต่างมีจุดเด่นในการปรับแต่งสำหรับสถานการณ์ภาษาจีน

\booksection{สำนัก Agent บนเทอร์มินัล}

Agent บนเทอร์มินัลคือสำนักที่โตเร็วที่สุดในปี 2026 แนวคิดหลัก: AI ไม่ควรเป็นแค่ปลั๊กอินของตัวแก้ไขโค้ด แต่ควรจัดการสภาพแวดล้อมการพัฒนาทั้งหมดได้เหมือนนักพัฒนา — อ่านเขียนไฟล์ รันคำสั่ง ดูผลลัพธ์ ติดตั้ง dependency Claude Code คือตัวแทนของสำนักนี้

\booksubsection{Claude Code}

Claude Code คือเครื่องมือเขียนโค้ดด้วย AI บนเทอร์มินัลจาก Anthropic รายได้เฉลี่ยต่อปีแตะระดับ 2.5 พันล้านดอลลาร์สหรัฐไปแล้วตั้งแต่กลางปี 2026 วิธีทำงานของมันต่างจาก IDE แบบดั้งเดิมโดยสิ้นเชิง: คุณบอกงานเป็นภาษาธรรมชาติในเทอร์มินัล แล้ว Claude Code จะอ่านไฟล์ที่เกี่ยวข้อง แก้โค้ด รันเทสต์ ดูผล และแก้ไขวนซ้ำด้วยตัวเอง — ตลอดกระบวนการคุณไม่ต้องเปิดตัวแก้ไขโค้ดเลย

จุดแข็งของ Claude Code คือ “การรับรู้ทั้งโปรเจกต์”: มันไม่ถูกจำกัดด้วยหน้าต่างของตัวแก้ไข เดินสำรวจไดเรกทอรีโปรเจกต์ทั้งหมดได้อิสระ และเข้าใจความสัมพันธ์ระหว่างไฟล์ มันยังรองรับไฟล์ความจำโปรเจกต์ CLAUDE.md ทำให้ AI รู้เบื้องหลังโปรเจกต์ตั้งแต่ต้นแต่ละเซสชัน สำหรับงานอย่าง refactor โปรเจกต์ใหญ่ แก้ไขเป็นชุด หรือเขียนเทสต์ Agent บนเทอร์มินัลทำได้มีประสิทธิภาพกว่าปลั๊กอิน IDE

\booksubsection{Codex CLI และ Aider}

Codex CLI ของ OpenAI ให้ประสบการณ์เขียนโค้ดบนเทอร์มินัลแบบคล้ายกัน และรองรับฟังก์ชัน /import — นำเข้าโค้ดเบสที่มีอยู่ให้ AI เข้าใจได้ Aider เป็นเครื่องมือเขียนโค้ดด้วย AI บนเทอร์มินัลแบบโอเพนซอร์ส รองรับหลายโมเดล (GPT, Claude, DeepSeek ฯลฯ) เหมาะกับผู้ใช้ที่อยากสลับโมเดลได้ยืดหยุ่น เครื่องมือเหล่านี้แม้จะมีจุดเด่นต่างกัน แต่แนวคิดหลักเหมือนกัน: ให้ AI มีความสามารถควบคุมสภาพแวดล้อมการพัฒนาแบบครบวงจรภายในเทอร์มินัล

\booksection{สำนัก Builder บนคลาวด์}

Builder บนคลาวด์คือสำนักที่เหมาะที่สุดกับผู้ใช้ที่เริ่มจากศูนย์ คุณไม่ต้องติดตั้งซอฟต์แวร์ใดเลย แค่เปิดเบราว์เซอร์ บอกความต้องการ แล้ว AI จะสร้างโปรเจกต์เต็มรูปแบบบนคลาวด์พร้อมพรีวิวให้ทันที

\begin{itemize}
  \item v0 (จาก Vercel): เก่งการสร้างคอมโพเนนต์ UI และหน้าเว็บฝั่ง front-end เหมาะกับทำต้นแบบเว็บอย่างรวดเร็ว
  \item Bolt: รองรับการสร้างแอป full-stack จากหนึ่งประโยค ครบทั้ง front-end, back-end และฐานข้อมูล
  \item Replit Agent: ทำงานใน IDE บนคลาวด์ของ Replit ไม่เพียงสร้างโปรเจกต์ แต่รันและ deploy บนคลาวด์ได้โดยตรง
\end{itemize}

จุดแข็งร่วมของเครื่องมือเหล่านี้คือ “เริ่มใช้ได้เลยไม่ต้องตั้งค่า” — เปิดเบราว์เซอร์ก็เริ่มได้ ข้อจำกัดคือ: เมื่อโปรเจกต์ซับซ้อนขึ้น สภาพแวดล้อมแก้ไขบนคลาวด์จะยืดหยุ่นไม่เท่า IDE ในเครื่อง ดังนั้นมันเหมาะที่สุดกับช่วงต้นแบบ “จาก 0 สู่ 1”: ใช้ Builder บนคลาวด์ตรวจแนวคิดอย่างรวดเร็ว พอผ่านแล้วค่อยย้ายมาลงลึกพัฒนาต่อบน IDE ในเครื่อง

\booksection{เครื่องมือเปลี่ยน วิธีคิดไม่เปลี่ยน}

ถ้าพรุ่งนี้เครื่องมือเขียนโค้ดด้วย AI หายไปทั้งหมด นิสัยสามข้อที่คุณควรยังคงไว้คือ: คิดให้ชัดก่อนลงมือ ทำทีละก้าวเล็ก ๆ และตรวจสอบทุกครั้งที่ทำเสร็จ เครื่องมือคือตัวขยาย วิธีการคือรากฐาน ตัวขยายเปลี่ยนได้ แต่รากฐานขยับไม่ได้

ย้อนดูการเปลี่ยนผ่านของเครื่องมือตั้งแต่ปี 2024 ถึง 2026: ปี 2024 GitHub Copilot เป็นจ้าวของวงการ ปี 2025 Cursor ขึ้นมาแทนที่ ปี 2026 Claude Code พุ่งทะยานก้าวขึ้นมา ทุกครั้งที่เครื่องมือใหม่เปิดตัว ก็จะมีคนประกาศว่า “ครั้งนี้ไม่เหมือนเดิม” แต่สุดท้ายพบว่า: คนที่ใช้เครื่องมือใหม่ได้ดี ล้วนเป็นพวกที่วางรากฐานวิธีคิดไว้แน่น พวกเขาเปลี่ยนเครื่องมือใช้เวลาปรับตัวแค่ไม่กี่วัน ส่วนคนที่เรียนแค่วิธีกดปุ่มของเครื่องมือ พอเปลี่ยนทีก็ต้องเริ่มเรียนใหม่ตั้งแต่ต้น

ดังนั้นเส้นทางการเรียนรู้ที่หนังสือเล่มนี้แนะนำคือ: เริ่มต้นด้วยเครื่องมือหนึ่งตัว (ตัวไหนก็ได้) แล้วซึมซับหลักคิดทั้งห้าจากการลงมือจริง เมื่อวิธีคิดของคุณแน่นพอแล้ว ค่อยลองเครื่องมือจากต่างสำนัก เพื่อสัมผัสความต่างของปรัชญาการออกแบบ สุดท้ายคุณจะได้ชุดเครื่องมือของตัวเอง — เช่น ใช้ Cursor พัฒนางานประจำวัน ใช้ Claude Code ทำ refactor กับเทสต์ และใช้ v0 ทำต้นแบบเร็ว ๆ ชุดเครื่องมือเป็นเรื่องของแต่ละคน แต่วิธีคิดคือของกลางที่ใช้ร่วมกันได้

\begin{tcbpractice}{ลงมือทำ: จัด "ชุดอาวุธประจำตัว" ให้ตัวเองสักชุด}
ใช้สามมุมมองการตัดสินใจ (พื้นฐานการเขียนโค้ด ช่วงของโปรเจกต์ ความต้องการ deploy) จัดตำแหน่งตัวเอง: ถ้าคุณเริ่มจากศูนย์ ลอง Bolt หรือ v0 ทำหน้าเว็บหนึ่งหน้า; มีพื้นฐาน ใช้ Cursor หรือ Trae ทำงานไล่เวอร์ชันหนึ่งงาน; อยากท้าทาย รัน Claude Code บนเทอร์มินัลสักครั้ง จากนั้นตอบตัวเองว่า: สำนักไหนตรงกับช่วงปัจจุบันของคุณที่สุด?

\end{tcbpractice}

\begin{tcbsummary}{สรุปบทนี้}
\begin{itemize}
  \item สามสำนัก: AI ในตัว IDE (พัฒนาประจำวัน), Agent บนเทอร์มินัล (งานระดับโปรเจกต์), Builder บนคลาวด์ (ต้นแบบสำหรับมือใหม่)
  \item เลือกเครื่องมือด้วยสามคำถาม: พื้นฐานการเขียนโค้ด? ช่วงของโปรเจกต์? ความต้องการ deploy?
  \item โหมด Agent ของ Cursor 3 แซงโหมดเติมโค้ด Tab เป็นครั้งแรก; รายได้เฉลี่ยต่อปีของ Claude Code แตะ 2.5 พันล้านดอลลาร์สหรัฐ
  \item Builder บนคลาวด์เหมาะกับช่วง "จาก 0 สู่ 1" ตรวจแนวคิดผ่านแล้วค่อยย้ายกลับมาลงลึกบน IDE ในเครื่อง
  \item เครื่องมือปรับรุ่นทุกสามเดือน: คนที่เรียนวิธีคิดเปลี่ยนเครื่องมือแล้วปรับตัวได้ในไม่กี่วัน
\end{itemize}

\end{tcbsummary}
